\paragraph{\texorpdfstring{$(4,2^5)$}{(4,2^5)} 分歧型的 Nielsen 计数、辫单值群与 Jacobian 三挠几何}
以下结论在本节既定 \(S_4\)-正则覆盖记号下成立，并固定一个阶 \(4\) 分歧点为“无穷远位”。
记 \(\mathcal C_4\subset S_4\) 为四循环共轭类，\(\mathcal C_2\subset S_4\) 为换位共轭类。

\begin{conclusion}[内 Nielsen 类基数 \(160\) 与 \(B_5\) 传递性]\label{con:xi-terminal-zm-s4-nielsen-in-160-b5-transitive}
设
\[
\mathrm{Ni}^{\mathrm{in}}(S_4;\mathcal C_4,\mathcal C_2^5)
:=\left\{
(g_\infty,g_1,\ldots,g_5)\in S_4^6\ \middle|\
\begin{array}{l}
g_\infty\in \mathcal C_4,\ g_i\in\mathcal C_2,\\
g_\infty g_1\cdots g_5=1,\\
\langle g_\infty,g_1,\ldots,g_5\rangle=S_4
\end{array}
\right\}\Big/\!\sim,
\]
其中 \(\sim\) 为同时共轭。
则
\[
\bigl|\mathrm{Ni}^{\mathrm{in}}(S_4;\mathcal C_4,\mathcal C_2^5)\bigr|=160.
\]
并且 \(5\)-股辫群 \(B_5\)（仅作用于 \((g_1,\ldots,g_5)\)）在该 \(160\) 元集合上作用传递。
\end{conclusion}
\begin{proof}
设 Hurwitz 生成元 \(\sigma_i\ (1\le i\le 4)\) 作用为
\[
(g_i,g_{i+1})\longmapsto \bigl(g_{i+1},g_{i+1}^{-1}g_i g_{i+1}\bigr).
\]
脚本
\texttt{scripts/exp\_fold\_zm\_s4\_s3\_nielsen\_monodromy\_audit.py}
对 \(S_4\) 完全枚举给出原始满足约束的六元组数为 \(3840\)，按同时共轭商去后得到
\[
3840/|S_4|=3840/24=160.
\]
同一脚本对 \(\sigma_i,\sigma_i^{-1}\) 诱导图执行轨道遍历，证得唯一轨道，故 \(B_5\) 作用传递。证毕。
\end{proof}

\begin{conclusion}[\texorpdfstring{$S_4\twoheadrightarrow S_3$}{S4 to S3} 诱导的 \(40\) 块四重分块]\label{con:xi-terminal-zm-s4-to-s3-nielsen-40-blocks-of-4}
令 \(V_4^{\mathrm{norm}}\triangleleft S_4\) 为 Klein 四元正规子群，并取商映射
\[
\pi:S_4\longtwoheadrightarrow S_4/V_4^{\mathrm{norm}}\cong S_3.
\]
设 \(\mathcal T\subset S_3\) 为换位共轭类，则 \(\pi\) 诱导映射
\[
\pi_*:\mathrm{Ni}^{\mathrm{in}}(S_4;\mathcal C_4,\mathcal C_2^5)\longrightarrow
\mathrm{Ni}^{\mathrm{in}}(S_3;\mathcal T^6).
\]
并满足：
\begin{enumerate}
  \item \(\bigl|\mathrm{Ni}^{\mathrm{in}}(S_3;\mathcal T^6)\bigr|=40\)；
  \item \(\pi_*\) 满射，且每个纤维基数恒为 \(4\)；
  \item 上述 \(160\) 元集合因此分解为 \(40\) 个 \(B_5\)-不变块，每块大小 \(4\)。
\end{enumerate}
\end{conclusion}
\begin{proof}
同一脚本在 \(S_3\) 侧完全枚举得原始六元组总数 \(240\)，按同时共轭商去得到
\[
240/|S_3|=240/6=40.
\]
对每个 \(S_4\)-内类代表施行 \(\pi\) 并计数纤维，得到纤维大小分布恒为 \(4\)。
由 Hurwitz 作用与群同态、同时共轭均可交换，\(\pi_*\) 的纤维自动组成 \(B_5\)-不变分块系统。证毕。
\end{proof}

\begin{conclusion}[块上辫单值群识别为 \(\mathrm{PSp}_4(\mathbf F_3)\)]\label{con:xi-terminal-zm-braid-monodromy-psp4f3-identification}
记 \(\Gamma_{S_3}\subset S_{40}\) 为 \(B_5\) 在
\(\mathrm{Ni}^{\mathrm{in}}(S_3;\mathcal T^6)\) 上的像，\(\Gamma_{S_4}\subset S_{160}\) 为 \(B_5\) 在
\(\mathrm{Ni}^{\mathrm{in}}(S_4;\mathcal C_4,\mathcal C_2^5)\) 上的像。
则：
\begin{enumerate}
  \item \(|\Gamma_{S_3}|=25920\)，且 \(\Gamma_{S_3}\) 中心平凡、导群等于自身；
  \item 点稳定子次轨道分解为 \((1,12,27)\)；
  \item 因而 \(\Gamma_{S_3}\cong \mathrm{PSp}_4(\mathbf F_3)\)；
  \item \(160\) 点作用为带 \(40\) 个 \(4\)-块的非本原作用，其块上商像即 \(\Gamma_{S_3}\)。
\end{enumerate}
进一步地，若记
\[
|\Gamma_{S_4}|=2^{102}\cdot 3^{44}\cdot 5,
\]
则块核满足
\[
|K|=\frac{|\Gamma_{S_4}|}{|\Gamma_{S_3}|}=2^{96}\cdot 3^{40}.
\]
\end{conclusion}
\begin{proof}
脚本证书给出
\[
|\Gamma_{S_3}|=25920,\quad Z(\Gamma_{S_3})=1,\quad [\Gamma_{S_3},\Gamma_{S_3}]=\Gamma_{S_3},
\]
并给出次轨道 \((1,12,27)\)。
另一方面，\(\mathrm{PSp}_4(\mathbf F_3)\) 的阶为 \(25920\)，其在 \(\mathbf P^3(\mathbf F_3)\) 的 \(40\) 点作用为秩 \(3\)，次轨道正为 \((1,12,27)\)。
由这些群论不变量可识别块上单值群为 \(\mathrm{PSp}_4(\mathbf F_3)\)。

块结构由结论 \ref{con:xi-terminal-zm-s4-to-s3-nielsen-40-blocks-of-4} 给出，故 \(160\) 点作用非本原且块上商像为 \(\Gamma_{S_3}\)。
关于 \(|\Gamma_{S_4}|\) 与 \(|K|\) 的数值，脚本中保留了对应 Schreier--Sims 证书值及
\[
|\Gamma_{S_4}|/|\Gamma_{S_3}|=2^{96}3^{40}
\]
的一致性核验。证毕。
\end{proof}

\begin{conclusion}[四十个分块与属 \(2\) Jacobian 的循环 \(3\)-挠子群]\label{con:xi-terminal-zm-40-blocks-equal-cyclic-3torsion-subgroups}
设
\[
\varphi:X\to\PP^1
\]
为本节固定的 \(S_4\)-Galois 覆盖，分歧型 \((4,2,2,2,2,2)\)。
记
\[
C:=X/A_4,\qquad C_6:=X/V_4^{\mathrm{norm}}.
\]
则：
\begin{enumerate}
  \item \(C\to\PP^1\) 为在六个分歧值处分歧的二次覆盖，故 \(g(C)=2\)；
  \item \(C_6\to C\) 为次数 \(3\) 的循环无分歧覆盖；
  \item 该覆盖由 \(L\in \mathrm{Pic}^0(C)[3]\setminus\{0\}\)（等价于循环子群 \(\langle L\rangle\)）确定；
  \item 因 \(\mathrm{Pic}^0(C)[3]\cong (\ZZ/3\ZZ)^4\)，非平凡循环子群个数为
  \[
  \#\mathbf P^3(\mathbf F_3)=\frac{3^4-1}{3-1}=40;
  \]
  \item 结论 \ref{con:xi-terminal-zm-s4-to-s3-nielsen-40-blocks-of-4} 的 \(40\) 个块可几何识别为这些循环 \(3\)-挠子群。
\end{enumerate}
\end{conclusion}
\begin{proof}
第 \(1\) 点由 \(S_4/A_4\cong C_2\) 及结论
\ref{con:xi-terminal-zm-s4-normal-subgroup-quotients-explicit-models-genus}
直接得到。
第 \(2\) 点由 \(A_4/V_4^{\mathrm{norm}}\cong C_3\) 与结论
\ref{con:xi-terminal-zm-c6-over-c-etale-c3-pic3-class}
给出的无分歧性得到。
第 \(3\) 点为无分歧循环 \(3\)-覆盖与 \(3\)-挠线丛的标准对应。
第 \(4\) 点为有限线性代数计数。
第 \(5\) 点由结论 \ref{con:xi-terminal-zm-s4-to-s3-nielsen-40-blocks-of-4} 的基数与纤维分块，与上述 \(40\)-点参数空间同基数同构识别给出。证毕。
\end{proof}

\begin{theorem}[块轨道的辛正交刻画与强正则结构]\label{thm:xi-terminal-zm-block-orbits-symplectic-orthogonality-srg}
在结论 \ref{con:xi-terminal-zm-40-blocks-equal-cyclic-3torsion-subgroups} 的识别下，将
\[
\mathcal B:=\mathrm{Ni}^{\mathrm{in}}(S_3;\mathcal T^6)
\]
视为 \(V:=\mathrm{Pic}^0(C)[3]\cong (\ZZ/3\ZZ)^4\) 的一维 \(\FF_3\)-子空间（即循环子群）之集合。
记主极化所诱导的 Weil 配对为 \(e_3:V\times V\to\mu_3\)，并对两条线 \(\ell_1,\ell_2\subset V\) 定义关系
\[
\ell_1\perp \ell_2
\quad:\Longleftrightarrow\quad
e_3(x,y)=1\ \text{对任意}\ x\in\ell_1\setminus\{0\},\ y\in\ell_2\setminus\{0\}.
\]
则：
\begin{enumerate}
  \item 在 \(\mathcal B\) 上存在唯一的 \(\Gamma_{S_3}\)-不变非平凡对称二元关系，其在点稳定子下的两条非平凡子轨道基数为 \((12,27)\)；并且该关系与 \(\perp\) 完全一致。
  \item 因而以 \(\mathcal B\) 为顶点、以 \(\ell\neq\ell'\) 且 \(\ell\perp\ell'\) 为边的图为强正则图，参数为
  \[
  (v,k,\lambda,\mu)=(40,12,2,4).
  \]
  \item 其复置换表示分解满足
  \[
  \CC[\mathcal B]\cong \mathbf 1\oplus W_{24}\oplus W_{15},
  \qquad
  \dim W_{24}=24,\ \dim W_{15}=15,
  \]
  等价地邻接算子谱为 \(12,2,-4\)（重数分别为 \(1,24,15\)）。
\end{enumerate}
\end{theorem}

\begin{proof}
由结论 \ref{con:xi-terminal-zm-braid-monodromy-psp4f3-identification}，\(\Gamma_{S_3}\cong \mathrm{PSp}_4(\mathbf F_3)\) 在 \(\mathcal B\) 上为秩 \(3\) 的传递作用，点稳定子次轨道分解为 \((1,12,27)\)；因此任一 \(\Gamma_{S_3}\)-不变对称二元关系由两条非平凡子轨道之一唯一确定，从而至多存在一个非平凡关系。
\par
另一方面，在四维辛空间 \(V\cong\FF_3^4\) 的射影点集 \(\PP(V)\) 上，固定一条线 \(\ell\) 后其正交补 \(\ell^\perp\) 为三维子空间，故 \(\PP(\ell^\perp)\) 含
\(
(3^3-1)/(3-1)=13
\)
个点；去除 \(\ell\) 自身尚余 \(12\) 点，正给出 \(\ell\) 的正交邻域
\(
\mathcal B_\perp(\ell):=\{\ell'\neq\ell:\ \ell'\subset\ell^\perp\}
\)。
余下 \(27\) 点为非正交类。
因此 \(\perp\) 给出一条 \((12,27)\) 子轨道划分，并与上述唯一性一致，得到第 (1) 条。
\par
对强正则参数，若 \(\ell_1\perp\ell_2\) 且 \(\ell_1\neq\ell_2\)，则 \(\ell_1^\perp\cap \ell_2^\perp\) 为二维子空间，其射影点数为 \(4\)；其中包含 \(\ell_1,\ell_2\)，故公共邻点数为 \(2\)，即 \(\lambda=2\)。
若 \(\ell_1\not\perp\ell_2\)，则 \(\ell_1,\ell_2\not\subset \ell_1^\perp\cap\ell_2^\perp\)，公共邻点数为 \(4\)，即 \(\mu=4\)。
于是得到第 (2) 条。
第 (3) 条为强正则图的一般谱结论，与 Klingen Hecke 分解中的同一谱值一致（参见定理 \ref{thm:conclusion-m2-level3-delta0-inertia-klingen-charpoly}）。证毕。
\end{proof}

\begin{proposition}[由强正则图内生重构辛极空间 \(W(3)\) 的线系]\label{prop:xi-terminal-zm-reconstruct-W3-from-srg}
设图为定理 \ref{thm:xi-terminal-zm-block-orbits-symplectic-orthogonality-srg} 的强正则图。
则其所有极大团的大小均为 \(4\)。
记 \(\mathcal L\) 为全部 \(4\)-团之集合，则：
\begin{enumerate}
  \item \(|\mathcal L|=40\)；
  \item 每个顶点恰落在 \(4\) 个 \(4\)-团中；
  \item \((\mathcal B,\mathcal L,\in)\) 构成阶 \(3\) 的广义四边形（辛极空间）\(W(3)\)：每点在 \(4\) 条线上、每线含 \(4\) 点，且满足广义四边形公理。
\end{enumerate}
\end{proposition}

\begin{proof}
在 \(\mathcal B\simeq \PP(V)\) 的识别下，\(4\)-团等价于 \(\PP(M)\)，其中 \(M\subset V\) 为二维全各向同性子空间（Lagrangian 平面）。
确有 \(|\PP(M)|=(3^2-1)/(3-1)=4\)，且任意两点正交，故给出 \(4\)-团。
反向地，若四点两两正交，则其张成的子空间必为二维各向同性子空间，从而团必为某个 \(\PP(M)\)。
\par
计数：固定一点 \(\ell\in\PP(V)\)，包含 \(\ell\) 的 Lagrangian 平面个数等于 \(\PP(\ell^\perp/\ell)\) 的点数；而 \(\ell^\perp/\ell\) 为二维 \(\FF_3\)-空间，故该点数为 \((3^2-1)/(3-1)=4\)，得到第 (2) 条。
再由总旗数计数得 \(|\mathcal L|\cdot 4=|\mathcal B|\cdot 4=160\)，故 \(|\mathcal L|=40\)，得第 (1) 条。
广义四边形公理为辛极空间的标准性质：对点 \(\ell\) 与不含 \(\ell\) 的 Lagrangian 平面 \(M\)，存在唯一 \(\ell'\in\PP(M)\) 使 \(\ell\perp\ell'\)，从而唯一线 \(\langle \ell,\ell'\rangle\) 与 \(M\) 相交于 \(\ell'\)。证毕。
\end{proof}

\begin{theorem}[\texorpdfstring{\(\pi_*\)}{pi*} 的四重纤维与“通过一点的四条线”]\label{thm:xi-terminal-zm-pistar-fiber-equals-four-lines-through-point}
\leanverified{paper\_xi\_terminal\_zm\_pistar\_fiber\_equals\_four\_lines\_through\_point}
保持结论 \ref{con:xi-terminal-zm-s4-to-s3-nielsen-40-blocks-of-4} 的四重覆盖
\(
\pi_*:\mathrm{Ni}^{\mathrm{in}}(S_4;\mathcal C_4,\mathcal C_2^5)\to \mathcal B
\)。
固定 \(b\in\mathcal B\)，令 \(B:=\pi_*^{-1}(b)\) 为其纤维。
在定理 \ref{thm:xi-terminal-zm-block-orbits-symplectic-orthogonality-srg} 与命题 \ref{prop:xi-terminal-zm-reconstruct-W3-from-srg} 的识别下，将 \(b\) 视为 \(V\) 中一条线 \(\ell\subset V\)，并记
\[
\mathcal L(\ell):=\{M\in\mathcal L:\ \ell\in M\},
\qquad |\mathcal L(\ell)|=4.
\]
则存在一个由单值作用所确定的同构类，使得纤维 \(B\) 的点稳定子单值表示与 \(\mathcal L(\ell)\) 的自然作用同构；从而在该同构类意义下可将 \(B\) 识别为“通过点 \(\ell\) 的四条线”集合 \(\mathcal L(\ell)\)。
在此识别下，\(\ell\) 的正交邻域分裂为四个互不交叠的三元组
\[
\mathcal B_\perp(\ell)
=
\bigsqcup_{M\in\mathcal L(\ell)}\bigl(\PP(M)\setminus\{\ell\}\bigr),
\qquad |\mathcal B_\perp(\ell)|=12,
\]
并且上述分裂由纤维 \(B\) 内在标记。
\end{theorem}

\begin{proof}
由结论 \ref{con:xi-terminal-zm-braid-monodromy-psp4f3-identification}，块上单值群为 \(G=\Gamma_{S_3}\cong \mathrm{PSp}_4(\mathbf F_3)\)，并且 \(G\) 在 \(\mathcal B\) 上秩为 \(3\)。
设 \(P=\mathrm{Stab}_G(\ell)\)。
在辛几何中，集合 \(\mathcal L(\ell)\) 与 \(\PP(\ell^\perp/\ell)\cong \PP^1(\FF_3)\) 等价，且 \(P\) 在其上诱导的作用同构于 \(\mathrm{PGL}_2(\FF_3)\cong S_4\) 的自然四点作用，因此该四点传递作用在同构意义下唯一。
另一方面，\(\pi_*\) 给出大小为 \(4\) 的纤维，且纤维上的单值表示由点稳定子诱导；于是纤维单值表示与 \(\mathcal L(\ell)\) 的自然表示在同构类意义下必须一致。
最后的三元组分裂由 \(\PP(M)\) 各含四点且不同 \(M\) 仅交于 \(\ell\) 直接推出。证毕。
\end{proof}

\begin{corollary}[\texorpdfstring{\((3,3)\)}{(3,3)}-同源核的四重扩充律]\label{cor:xi-terminal-zm-cyclic3-to-33-isogeny-four-extensions}
令 \(A=\mathrm{Pic}^0(C)\) 为主极化阿贝尔曲面。
对任一非平凡循环子群 \(\ell\subset A[3]\)，包含 \(\ell\) 的极大各向同性子群 \(M\subset A[3]\) 恰有 \(4\) 个，并与 \(\mathcal L(\ell)\) 等势。
对每个 \(M\) 商 \(A/M\) 继承主极化，从而诱导一条 \((3,3)\)-同源
\[
\varphi_M:A\to A/M.
\]
因此，“给定循环 \(3\)-同源核 \(\ell\)”具有且仅具有四种方式可扩充为一个 \((3,3)\)-同源核 \(M\)。
\end{corollary}

\begin{proof}
极大各向同性二维子空间与 \((3,3)\)-同源核的对应由点--线四重纤维识别给出（见定理 \ref{thm:xi-terminal-zm-pistar-fiber-equals-four-lines-through-point}）。
而包含给定 \(\ell\) 的极大各向同性平面数目为 \(|\PP(\ell^\perp/\ell)|=4\)，与命题 \ref{prop:xi-terminal-zm-reconstruct-W3-from-srg} 一致。证毕。
\end{proof}

\begin{proposition}[线系的强正则对偶与旗空间的双 \(4\)-重投影]\label{prop:xi-terminal-zm-line-srg-and-flag-double-projection}
令 \(\mathcal L\) 为命题 \ref{prop:xi-terminal-zm-reconstruct-W3-from-srg} 的线系，并在 \(\mathcal L\) 上定义邻接关系
\[
M\sim N\quad:\Longleftrightarrow\quad \card{M\cap N}=3
\quad(\text{即 } M\cap N\ \text{为一条线}).
\]
则：
\begin{enumerate}
  \item \((\mathcal L,\sim)\) 亦为参数 \((40,12,2,4)\) 的强正则图；
  \item 点–线关联图 \(\mathcal I\subset \mathcal B\times\mathcal L\)（\((\ell,M)\in\mathcal I\iff \ell\subset M\)）为 \((4,4)\)-正则二部图，边数为 \(160\)；
  \item \(\mathcal I\to\mathcal B\) 与 \(\mathcal I\to\mathcal L\) 均为四重投影，并在组合意义下互为对偶：每个点邻域的 \(12\) 个正交邻点被唯一分裂为四个三元组，对应于通过该点的四条线；线邻域亦同理。
\end{enumerate}
\end{proposition}

\begin{proof}
在辛极空间 \(W(3)\) 中点与线完全对偶。
固定 \(M\in\mathcal L\)，其包含四条线 \(\ell\subset M\)；每条 \(\ell\) 被恰好四条线包含，其中之一为 \(M\) 本身，故与 \(M\) 交于 \(\ell\) 的其他线有三条。
因此 \(M\) 的邻接度为 \(4\cdot 3=12\)，余下为 \(27\)，与点图同参。
其余参数由同样的二维交与公共邻点计数得到，与点图的 \((\lambda,\mu)=(2,4)\) 完全平行。
二部图与双四重投影为命题 \ref{prop:xi-terminal-zm-reconstruct-W3-from-srg} 的计数结论之直接改写；对偶分裂由
\(
\mathcal B_\perp(\ell)=\bigsqcup_{M\in\mathcal L(\ell)}(\PP(M)\setminus\{\ell\})
\)
给出。证毕。
\end{proof}

\begin{conjecture}[内 Hurwitz 空间的抛物型 \(3\)-水平结构模解释]\label{conj:xi-terminal-zm-hurwitz-parabolic-level3-moduli}
设 \(\mathcal H_{S_3}^{\mathrm{in}}(\mathcal T^6)\) 与 \(\mathcal H_{S_4}^{\mathrm{in}}(\mathcal C_4,\mathcal C_2^5)\) 分别为相应分歧型别的内 Hurwitz 空间连通分量。
则期望存在一条与结论 \ref{con:xi-terminal-zm-40-blocks-equal-cyclic-3torsion-subgroups} 的 Jacobian 三挠识别与定理 \ref{thm:xi-terminal-zm-pistar-fiber-equals-four-lines-through-point} 的点–线–旗几何同时相容的双有理同构链，使得：
\begin{enumerate}
  \item \(\mathcal H_{S_3}^{\mathrm{in}}(\mathcal T^6)\) 与某个抛物型 \(3\)-水平结构的 Siegel 模叠（等价地：主极化阿贝尔曲面携带一条 \(3\)-挠线 \(\ell\subset A[3]\) 的模叠）在双有理意义下同构；
  \item 覆盖 \(\mathcal H_{S_4}^{\mathrm{in}}(\mathcal C_4,\mathcal C_2^5)\to \mathcal H_{S_3}^{\mathrm{in}}(\mathcal T^6)\) 对应于进一步选择一条包含 \(\ell\) 的极大各向同性子群 \(M\subset A[3]\)（即选择 \((3,3)\)-同源核），其几何纤维与 \(\PP(\ell^\perp/\ell)\cong \PP^1(\FF_3)\) 的四点截面一致。
\end{enumerate}
该猜想的可核验截面由 \(\Gamma_{S_3}\cong \mathrm{PSp}_4(\mathbf F_3)\) 的块上单值化、强正则结构与旗空间双投影给出；其几何提升需在代数簇层面构造相应的 \(A[3]\) 局部系统并验证其与 Hurwitz 单值数据的自然性一致。
\end{conjecture}

\begin{corollary}[Hurwitz 模空间不可约性与“属 \(2\)+\(3\)-挠数据”解释]\label{cor:xi-terminal-zm-hurwitz-irreducible-and-genus2-jac3torsion-moduli}
记 \(\mathcal H_{S_4}\) 为上述 \((4,2^5)\) 型 \(S_4\)-覆盖的 Hurwitz 模空间，
\(\mathcal H_{S_3}\) 为其 \(S_3\)-商型空间。则：
\begin{enumerate}
  \item \(\mathcal H_{S_4}\) 与 \(\mathcal H_{S_3}\) 均几何不可约，维数均为 \(3\)；
  \item 存在有限 \'{e}tale 映射 \(\mathcal H_{S_4}\to \mathcal H_{S_3}\)，泛纤维基数为 \(4\)；
  \item \(\mathcal H_{S_3}\) 可解释为“带六个分歧标记的属 \(2\) 曲线 \(C\) 及其 \(J(C)[3]\) 非平凡循环子群”之粗模空间。
\end{enumerate}
\end{corollary}
\begin{proof}
由结论 \ref{con:xi-terminal-zm-s4-nielsen-in-160-b5-transitive}、
\ref{con:xi-terminal-zm-s4-to-s3-nielsen-40-blocks-of-4} 的辫群传递性，Hurwitz 覆盖在分歧点配置空间上连通，从而几何不可约。
分歧点总数为 \(6\)，故维数为 \(6-3=3\)。
第 \(2\) 点由纤维大小恒为 \(4\) 得到。
第 \(3\) 点由结论 \ref{con:xi-terminal-zm-40-blocks-equal-cyclic-3torsion-subgroups} 给出的几何识别得到。证毕。
\end{proof}

\begin{proposition}[Hurwitz 生成元循环型的离散证书]\label{prop:xi-terminal-zm-hurwitz-generator-cycle-types-on-160-and-40}
在上述置换表示中，对任一标准生成元 \(\sigma_i\ (1\le i\le 4)\)：
\begin{enumerate}
  \item 在 \(160\) 点集合 \(\mathrm{Ni}^{\mathrm{in}}(S_4;\mathcal C_4,\mathcal C_2^5)\) 上，
  \[
  \sigma_i\sim 3^{36}2^{12}1^{28},
  \]
  因而置换阶为 \(6\)；
  \item 在 \(40\) 点集合 \(\mathrm{Ni}^{\mathrm{in}}(S_3;\mathcal T^6)\) 上，
  \[
  \sigma_i\sim 3^{9}1^{13},
  \]
  因而置换阶为 \(3\)。
\end{enumerate}
\end{proposition}
\begin{proof}
由脚本对四个 Hurwitz 生成元的循环分解直接给出：
\[
\{1^{28},2^{12},3^{36}\}\quad\text{与}\quad \{1^{13},3^{9}\},
\]
且四个生成元同型。置换阶由循环长度最小公倍数得到。证毕。
\end{proof}

\begin{conclusion}[模空间像的不可约性与一般点自同构群]\label{con:xi-terminal-zm-m16-locus-irreducible-generic-aut-s4}
令 \(\mathcal L\subset \mathcal M_{16}\) 为该 Hurwitz 家族在属 \(16\) 模空间中的像。
则：
\begin{enumerate}
  \item \(\mathcal L\) 为维数 \(3\) 的不可约代数簇；
  \item \(\mathcal L\) 一般点对应曲线 \(X\) 满足
  \[
  \mathrm{Aut}(X)=S_4.
  \]
\end{enumerate}
\end{conclusion}
\begin{proof}
第 \(1\) 点由推论
\ref{cor:xi-terminal-zm-hurwitz-irreducible-and-genus2-jac3torsion-moduli}
及模映射像的不可约性得到。
第 \(2\) 点中，对一般分歧配置，基底 \(\PP^1\) 上保持分歧集合的自同构群为平凡。
任何 \(X\) 的自同构都必须保持 \(S_4\)-覆盖结构并在基底诱导该平凡自同构，故仅剩甲板变换群 \(S_4\) 本身。证毕。
\end{proof}

\begin{remark}[可复现证书路径]\label{rem:xi-terminal-zm-s4-s3-nielsen-monodromy-certificate-path}
本组结论的离散计数与置换群证书由脚本
\texttt{scripts/exp\_fold\_zm\_s4\_s3\_nielsen\_monodromy\_audit.py}
生成，输出文件为
\texttt{artifacts/export/fold\_zm\_s4\_s3\_nielsen\_monodromy\_audit.json}。
\end{remark}

\endinput
